\paragraph{Independent audit of the fourteen shell bands and their lifts.}
Set $x=n/p$, $y=k/p$, and $\varepsilon=\lfloor2x-y\rfloor$.
The exact numerator high-digit root intervals are
\[
 [\lfloor ux-y\rfloor+1,\lfloor(u+10)x-y\rfloor],\qquad u=0,1,2,
\]
and the denominator intervals are
\[
 [\lceil-bx-y\rceil,\lfloor-ax-y\rfloor],
 \qquad(a,b)=(1,17),(2,16).
\]
These formulas include every integer endpoint correctly. Their changes occur
on lines $y=mx+c$, with
$m\in\{0,1,2,10,11,12,-1,-2,-16,-17\}$ and $c\in\mathbb Z$.
All intersections have $x$-denominator at most $29$. Thus the Farey-$29$
partition, together with every $y$-boundary and every open $y$-strip,
is an exhaustive finite arrangement certificate. For $p>29$ and
$0<\{n/p\}<1$, none of the omitted $x$-walls can occur.
The case $\{n/p\}=0$ is outside the selected bands.

The independently implemented certificate checks 270 $x$-cells and
5400 open or boundary strata. Exactly 47 selected cells merge into the
fourteen stated bands. Throughout every selected cell, $c(k)-d(k)\ge-1$;
on its minimal support $c(k)-d(k)=-1$, all ten tagged interval endpoints,
after subtraction of $\varepsilon$, are independent of $y$.
Keeping the interval tags is essential for the following uniform lift.

Write $x=t+f$, $t\in\mathbb Z_{\ge0}$, $0<f<1$. Then
$\varepsilon(t+f,y)=2t+\varepsilon(f,y)$. After subtracting $\varepsilon$,
the numerator endpoints change by $(u-2)t,(u+8)t$, and the denominator
endpoints change by $-(b+2)t,-(a+2)t$. These changes are independent of
$y$. Each numerator interval gains $10t$ roots; the denominator intervals
gain $16t$ and $14t$. Hence $c-d$ is unchanged, and the entire normalized
high-digit rational kernel remains independent of $k$ on minimal support,
for every integer $t$, not merely the finitely many tested lifts.

Assume $p>29$ and $p^2>30n$. Distinct high-digit roots differ by less than
$29n/p<p$ and hence by a $p$-adic unit. Also every harmonic high index is
at most $19n/p<p$. These inequalities justify local partial-fraction
extraction and harmonic reduction without any unrecorded divisibility loss.
The weaker cutoff $p^2>60n$ also suffices but is unnecessary.

As in the broad-shell proof, the multiplicity inequality gives
\[
 \mathcal Q(X)=(X^p-X)\overline{\mathcal N}(X)/
                       \overline{\mathcal D}(X)
 \in\mathbb F_p[X],\qquad\deg\mathcal Q=p-2.
\]
The polynomial vanishes outside minimal support. On minimal support the
local rational function $pF_n(k+pV)$ is $-\mathcal Q(k)$ times the
high-digit kernel. Translation $V\mapsto V+\varepsilon$ puts its poles
at the normalized endpoints just certified, and the relevant harmonic
indices become their negatives. Its harmonic readout is consequently a
constant depending on the $x$-cell and $t$, but independent of $k$.
Thus $p^{r-1}P_{r,n}$ is that constant times
$\sum_{k\in\mathbb F_p}\mathcal Q(k)=0$, for $r=3,4$.
In particular
\[
 v_p(Q_n)\ge1,\qquad v_p(P_{r,n})\ge2-r.
\]
Since $p\le19n$ in these bands and every lift,
the normalized pair $6C_nd_{19n}^{r}(Q_n,P_{r,n})$ has common valuation
at least $c(f)+2$, where
$c(f)=2\lfloor15f\rfloor-3\lfloor10f\rfloor$.
The certificate separately verifies that this equals the baseline
valuation plus one, both in the first period and in every later period.
There is therefore no double counting of the improvement.

The first-period density is exactly $195197/55440$. Including all integer
lifts gives the positive series
\[
 \sum_{(a,b)\in\mathcal B}\sum_{t=0}^{\infty}
 \left(\frac1{t+a}-\frac1{t+b}\right),
\]
where $\mathcal B$ is the fourteen-band set in the $x$ variable.
The first 1000 periods already contribute more than $3.674728195806$.
The finitely many fixed primes and primes $p^2\le30n$ contribute only
$o(n)$ to the logarithm of the divisor, so their exclusion does not change
the asymptotic density. This last assertion uses bounded profile exponents.
